\documentclass[a4paper,fleqn]{cas-sc}

\usepackage[numbers,sort&compress]{natbib}
\bibliographystyle{model1-num-names}
\usepackage{array}
\usepackage{tabularx}
\usepackage{graphicx}
\usepackage{xurl}
\graphicspath{{figures/}}

% Clickable, line-breakable repository URLs (hyperref loaded by cas-sc)
\newcommand{\githubrepo}{\url{https://github.com/subtlemyst/quaketwin-dhaka}}

% Declarations page: after Conclusion, before References
\newcommand{\declarationspage}{%
  \clearpage
  \begingroup
  \raggedright
  \sloppy
  \setlength{\parskip}{0.5\baselineskip}
  \section*{Data Availability}
  All input data are openly available: building footprints, road networks, health facilities, and lifeline infrastructure derive from OpenStreetMap (\copyright\ OpenStreetMap contributors, Open Database License) via the Geofabrik Bangladesh extract, and gridded population from the WorldPop 2020 constrained dataset (CC~BY~4.0). All derived layers (hazard grids, building risk profiles, diurnal exposure, response, and cascade outputs) are regenerated deterministically by the released pipeline.

  \section*{Code Availability}
  \phantomsection\label{sec:code_availability}
  The full pipeline is publicly available: OSM extraction (\texttt{scripts/extract\_dhaka\_osm.py}), hazard and enrichment stages (\texttt{scripts/run\_hazard\_scenario.py}, \texttt{scripts/build\_building\_profiles.py}), diurnal exposure and response models (\texttt{scripts/run\_diurnal\_exposure.py}, \texttt{scripts/run\_response\_model.py}), cascade simulation and physics-informed surrogate emulator (\texttt{scripts/run\_cascade.py}, \texttt{scripts/train\_cascade\_gnn.py}), publication scripts, configuration files (\texttt{config/}), and a README with step-by-step reproduction instructions.\par\medskip
  \noindent\textbf{GitHub}\par
  \githubrepo

  \section*{Ethics Statement}
  This study used only openly licensed geospatial and gridded population data; no human-subject data were collected or processed at individual level. Casualty figures are model-based scenario estimates for planning purposes.

  \section*{Declaration of Competing Interest}
  The author declares no known competing financial interests or personal relationships that could have influenced the work reported in this paper.
  \endgroup
  \clearpage
}

% Inline figures (CAS redefines floats; avoid caption package conflicts)
\makeatletter
\newcommand{\casinlinefig}[4][\linewidth]{%
  \par\vskip6pt
  \refstepcounter{figure}%
  {\sffamily\small\raggedleft
   \includegraphics[width=#1]{#2}\par\vskip6pt
   \textbf{\color{scolor}\figurename~\thefigure:}~#3\par}%
  \label{#4}%
  \par\vskip6pt
}
\makeatother

% References: dedicated final page(s); standard \twocolumn back matter (CAS-safe)
\makeatletter
\renewcommand{\bibsection}{%
  \clearpage
  \twocolumn[{\section*{References}\vspace{0.25\baselineskip}}]%
}
\AtBeginDocument{%
  \setlength{\bibsep}{0pt plus 0.2pt}%
}
\makeatother

\begin{document}
\let\WriteBookmarks\relax
\def\printorcid{}

\shorttitle{A Digital Twin of Earthquake Risk in Dhaka}
\shortauthors{Chowdhury}

\title[mode=title]{QuakeTwin: A Digital Twin of Earthquake Vulnerability,
Dynamic Population Exposure, and Cascading Lifeline Failure in Dhaka,
Bangladesh}

\author[1]{Md. Rafiul Azam Chowdhury}
\affiliation[1]{organization={Department of Computer Science and Engineering, United International University},
            addressline={},
            city={Dhaka},
            postcode={1212},
            country={Bangladesh}}

\begin{abstract}
Dhaka, one of the world's densest megacities, sits within 160~km of the Dauki
Fault, yet building-resolved earthquake risk assessments that account for
time-of-day population movement and interdependent lifeline failure remain
scarce. This study presents \textbf{QuakeTwin, an open and reproducible
digital twin of Greater Dhaka} that chains six coupled models: (i) a
GMPE logic-tree ground-motion model with Vs30 site amplification for a
moment magnitude ($M_w$)~7.2 Dauki scenario; (ii) fragility-based collapse
probabilities for 706{,}115 OpenStreetMap building footprints enriched with
WorldPop population under structural plausibility caps; (iii) a
population-conserving diurnal exposure model with HAZUS
severity-stratified casualties; (iv) an emergency-response model with
network shortest-path routing to 728 mapped health facilities; (v) a
300-iteration Monte Carlo cascade simulation over a 2{,}479-node lifeline
dependency graph, extended to a six-scenario ensemble spanning Dauki and
Madhupur faults; and (vi) a composite Earthquake Resilience Index (ERI)
over 487 building zones. Under the reference scenario, mean
soil-amplified peak ground acceleration reaches 0.22~g (mean collapse
probability 0.40), HAZUS-expected casualties at midday total 323{,}837, and
243 of 728 hospitals (33\%) exceed emergency capacity. Cascading failures
amplify lifeline loss far beyond direct seismic damage: substation failure
probability rises from 0.14 to 0.60 (4.3$\times$) and hospital functional
loss from 0.09 to 0.60 (7.1$\times$), with an expected rescue-delay factor
of 1.90. A nearer Madhupur $M_w$~7.0 event produces near-total lifeline
failure (100\% communication loss) and higher delay (2.34$\times$) despite
lower magnitude. A \textbf{physics-informed graph surrogate emulator}
trained on Monte Carlo cascade outputs generalizes across fault scenarios:
cross-scenario evaluation on held-out Madhupur graphs yields
$R^2 = 0.883$ (mean absolute error 0.078), versus $R^2 = 0.992$ on a
single-scenario node hold-out. Citywide ERI is 70.9 (very-high tier).
Outputs support \textbf{population-level preparedness planning and lifeline
hardening prioritization}, not building-specific engineering assessment.
The full pipeline is released as open source for data-scarce megacities.
\end{abstract}

\begin{keywords}
digital twin \sep earthquake risk \sep surrogate emulator \sep cascading
failure \sep dynamic exposure \sep Dhaka
\end{keywords}

{\hfuzz=120pt\maketitle}

\section{Introduction}
Dhaka concentrates more than 20~million people on Holocene floodplain
sediments at the collision boundary of the Indian and Eurasian plates.
Geodetic evidence indicates that the Indo-Burman subduction system east of
the city has accumulated sufficient strain to host large earthquakes
\cite{steckler2016}, and the Dauki Fault along the Shillong Plateau margin,
approximately 150--200~km north-east of the capital, is considered capable
of $M_w \geq 7$ events. Despite this, most published risk assessments for
Dhaka stop at ward-level hazard or vulnerability indices
\cite{cdmp2009,islam2020}: few resolve risk at the individual building, none
propagate that risk through time-of-day population movement, and none
quantify how interdependent lifeline systems --- electricity, communication,
road access, and health care --- fail together rather than in isolation.

These three gaps matter operationally. Building-level resolution determines
where retrofitting budgets go; diurnal exposure determines when an
earthquake of a given size is deadliest; and cascading lifeline failure
determines whether the emergency-response system that survives the shaking
can actually function. International experience is unambiguous on the last
point: in the 2011 T\=ohoku and 2023 Kahramanmara\c{s} earthquakes,
communication blackouts and hospital power loss degraded response far beyond
what structural damage alone predicted \cite{ouyang2014}. A \textbf{digital
twin} differs from a static geographic information system (GIS) overlay in
four ways \cite{ford2020,hamilton2024udt}: it maintains a persistent,
building-resolved representation of urban state; couples that state to
physics-based and graph-based simulators; supports repeatable what-if
scenario injection (magnitude, epicenter, retrofit); and is architected for
periodic or event-driven updating when live feeds (seismic instruments,
hospital occupancy, utility status) become available. QuakeTwin is therefore
a \textit{synthetic} twin today --- fully reproducible from open archives ---
with a surrogate layer that makes interactive exploration feasible before
real-time instrumentation is deployed.

This paper introduces \textbf{QuakeTwin}, a digital twin of Greater Dhaka
that addresses all three gaps in a single reproducible pipeline built
entirely from freely available data (OpenStreetMap \cite{osm2008}, WorldPop
\cite{worldpop}, and regional seismic-hazard literature). The contributions
are:
\begin{enumerate}
  \item a building-resolved risk layer for 706{,}115 footprints combining a
        weighted GMPE logic-tree, Vs30 site amplification, liquefaction
        susceptibility, and HAZUS-informed fragility curves \cite{hazus};
  \item a population-conserving diurnal exposure model with HAZUS
        severity-stratified casualties, redistributing rather than rescaling
        the WorldPop nighttime population across occupancy classes by hour;
  \item an emergency-response model with network shortest-path routing,
        hospital-overload, and rescue-priority estimates for 728 mapped
        health facilities;
  \item a six-scenario Monte Carlo cascade ensemble over an explicit
        lifeline dependency graph \cite{motter2002,ouyang2014}, quantifying
        fault-proximity effects (Dauki vs.\ Madhupur) on system-level failure;
  \item a \textbf{physics-informed surrogate emulator} --- implemented as
        GraphSAGE \cite{hamilton2017} after architecture ablation --- that
        accelerates Monte Carlo cascade simulation to millisecond latency
        without replacing the physics pipeline; and
  \item a composite Earthquake Resilience Index (ERI) integrating structural,
        emergency-capacity, lifeline, and accessibility components at the
        building-zone scale.
\end{enumerate}

\textbf{Epistemic caveat (read first).} The dominant uncertainty in this
study is not model architecture but \textit{data incompleteness}: 95.4\% of
building footprints lack OpenStreetMap construction tags and inherit a
single ``unknown'' fragility class. All downstream exposure, response, and
resilience metrics are conditional on this proxy layer and should be
interpreted as planning-order-of-magnitude estimates, not engineering
certificates for individual structures.

Table~\ref{tab:positioning} positions this study against regional assessments,
global hazard platforms, and digital-twin literature. QuakeTwin is
complementary to engines such as OpenQuake \cite{pagani2014} and the GEM
framework \cite{gem2023}: those systems excel at regional probabilistic
seismic hazard and portfolio loss; this work adds building-resolved open-data
stock, population-conserving diurnal exposure, network-routed response, and
explicit lifeline cascade simulation with a learned surrogate at city scale.

\begin{table}[pos=htbp]
  \centering
  \caption{Positioning against representative platforms and Dhaka studies. Bldg.\ = building-level; Diurnal = time-of-day exposure; Cascade = interdependent lifeline failure; Surrogate = learned simulation emulator; Open = fully open data and code.}
  \label{tab:positioning}
  \small
  \setlength{\tabcolsep}{2pt}
  \begin{tabularx}{\linewidth}{@{}>{\raggedright\arraybackslash}X*{5}{>{\centering\arraybackslash}p{0.09\linewidth}}@{}}
    \toprule
    Platform / study & Bldg. & Diurnal & Cascade & Surr. & Open \\
    \midrule
    CDMP \cite{cdmp2009} & No & No & No & No & No \\
    Islam et al.\ \cite{islam2020} & No & No & No & No & Part. \\
    HAZUS \cite{hazus} & Yes$^\dagger$ & No & No & No & No \\
    OpenQuake \cite{pagani2014} & Port. & No & No & No & Yes \\
    GEM \cite{gem2023} & Port. & No & No & No & Yes \\
    SimCenter workflows \cite{simcenter2020} & Yes & No & Part. & No & Yes \\
    ShakeMap \cite{wald2006} & No & No & No & No & Yes \\
    Urban DT reviews \cite{hamilton2024udt} & --- & No & No & No & --- \\
    \textbf{This study (QuakeTwin)} & \textbf{Yes} & \textbf{Yes} & \textbf{Yes} & \textbf{Yes} & \textbf{Yes} \\
    \bottomrule
  \end{tabularx}
  \par\smallskip
  {\footnotesize $^\dagger$Inventory-based, not OSM-resolved for Dhaka.}
\end{table}

Section~\ref{sec:method} formalizes the methodology and
Section~\ref{sec:experiments} reports empirical results for the
$M_w$~7.2 Dauki scenario.

\section{Proposed Method}
\label{sec:method}

The pipeline chains six stages --- hazard, building vulnerability, diurnal
exposure, emergency response, lifeline cascade with physics-informed surrogate
emulator, and
resilience indexing --- each consuming the previous stage's outputs and each
independently re-runnable from configuration files.

\subsection{Digital-Twin Architecture}
\label{sec:architecture}
QuakeTwin separates persistent \textit{state} (building and infrastructure
graphs), \textit{simulation} (hazard, fragility, cascade Monte Carlo), and
\textit{surrogate inference} (graph-based emulator). Figure~\ref{fig:dt_arch} shows the target
Phase~6 architecture: live feeds (seismic sensors, hospital status, utility
SCADA, OSM/WorldPop updates) refresh the state graph; the physics pipeline
calibrates cascade parameters; the surrogate emulator enables millisecond what-if
queries on the dashboard while full Monte Carlo runs remain available for
validation. This is distinct from a one-off GIS overlay or static HAZUS
portfolio run.

\casinlinefig[0.72\linewidth]{dt_architecture.pdf}{QuakeTwin digital-twin architecture: live data feeds, physics simulation chain, physics-informed surrogate emulator, and dashboard API.}{fig:dt_arch}

\subsection{Study Area and Data}
The study area covers Greater Dhaka (90.32--90.53$^{\circ}$E,
23.68--23.95$^{\circ}$N). All datasets are open (Table~\ref{tab:data}).
Only eight communication towers are mapped in OpenStreetMap for the study
area; communication-layer results are therefore indicative and flagged as a
data limitation throughout (Section~\ref{sec:limitations}).

\begin{table}[pos=htbp]
  \centering
  \caption{Data sources.}
  \label{tab:data}
  \small
  \setlength{\tabcolsep}{4pt}
  \begin{tabularx}{\linewidth}{@{}>{\raggedright\arraybackslash}p{0.34\linewidth}>{\raggedright\arraybackslash}X>{\raggedright\arraybackslash}p{0.26\linewidth}@{}}
    \toprule
    Layer & Source & Count / resolution \\
    \midrule
    Building footprints & OpenStreetMap (Geofabrik) & 706{,}115 polygons \\
    Road network & OpenStreetMap & highway-tagged ways \\
    Hospitals and clinics & OpenStreetMap & 728 facilities \\
    Power nodes & OpenStreetMap \texttt{power=*} & 61 \\
    Communication towers & OpenStreetMap & 8 \\
    Highway bridges & OpenStreetMap & 1{,}195 \\
    Population & WorldPop 2020 constrained \cite{worldpop} & $\sim$100~m grid \\
    Soil zones & Literature-informed proxy & 3 classes \\
    \bottomrule
  \end{tabularx}
\end{table}

\subsection{Scenario and Ground Motion}
\label{sec:hazard}
The reference scenario is a $M_w$~7.2 rupture on the Dauki Fault
(epicentral distance to central Dhaka $\approx 158$~km, depth 15~km).
Median bedrock peak ground acceleration (PGA) follows a compact
distance-saturating attenuation of Boore--Atkinson form \cite{boore2008},
\begin{equation}
\ln \mathrm{PGA} = a + b\,M_w - c \ln\!\left(R^2 + h^2\right),
\end{equation}
implemented as a \textbf{weighted logic-tree} of two crustal GMPE branches
(55\% calibrated Boore-style, 45\% Atkinson--Boore 2006 scaled for eastern
Bangladesh). Weights were set by requiring the weighted-mean bedrock PGA at
Dhaka ($\approx$158~km) to fall within the 0.06--0.25~g envelope of regional
studies \cite{cdmp2009} while preserving the relative spread between branches;
a one-at-a-time sensitivity over 40/60 and 60/40 splits changes mean
amplified PGA by less than 4\% (Section~\ref{sec:sensitivity}); perturbing
Vs30 by $\pm$20\% changes mean collapse probability by up to 8.6\%.
Site response samples
a USGS Vs30 mosaic \cite{borcherdt1994} (180--375~m/s over the study area,
with Borcherdt-style short-period amplification $F_a = (V_{\mathrm{ref}}/V_{s30})^{m_a}$ where $m_a$ decreases with bedrock PGA), combined with three
literature-informed soil-proxy zones (amplification factors 1.3--2.1).
Liquefaction susceptibility on each grid cell follows a simplified
Iwasaki-style index \cite{iwasaki1982,kramer1996}:
\begin{equation}
L = \min\!\left[1,\; S_{\mathrm{vs30}} \cdot f_{\mathrm{shake}}(\mathrm{PGA_{amp}}) \cdot g_{\mathrm{gw}}(\mathrm{zone}) \right],
\end{equation}
where $S_{\mathrm{vs30}}$ is base susceptibility from Vs30 class,
$f_{\mathrm{shake}}$ ramps from 0.2 to 1.0 as PGA increases, and
$g_{\mathrm{gw}}$ is a groundwater proxy elevated in marsh and Holocene
zones. This index is a \textit{planning proxy} pending borehole calibration.
PGA converts to Modified Mercalli intensity (MMI) via Worden et al.\
\cite{worden2012}. Both prespecified validation checks pass: bedrock PGA
(model mean 0.161~g) falls within the 0.06--0.25~g literature envelope, and
mean MMI (8.08) within the 6.0--8.5 envelope for Greater Dhaka under a large
eastern-fault event.

\subsection{Building Vulnerability}
\label{sec:buildings}
Each footprint is enriched with construction type (normalized from OSM
\texttt{building=*} tags), height (from \texttt{building:levels} or class
defaults), occupancy class, nearest hazard-grid attributes, fault distance,
and population. Population is the product of WorldPop density and effective
footprint area, bounded by an effective-footprint cap
(20{,}000~m\textsuperscript{2}) and a structural capacity cap of
$\mathrm{floors} \times \mathrm{area} / 10$~m\textsuperscript{2} per person;
these caps remove implausible single-building estimates (maximum 5{,}333
occupants after capping) while changing the citywide total by only $-2.8\%$
(final total 4.53~million).

Collapse probability follows a HAZUS-informed logistic fragility
\cite{hazus} on soil-amplified PGA,
\begin{equation}
P_{\mathrm{collapse}} = \min\!\left[1,\;
\frac{1}{1+e^{-k (\mathrm{PGA_{amp}} - \theta_c)}} + 0.25\,L \right],
\end{equation}
where $\theta_c$ is the construction-class median capacity (reduced for
taller buildings), $k$ the logistic slope, and $L \in [0,1]$ the
liquefaction index (boost capped at 0.20). A gradient-boosting model
trained on the same features is retained solely as a sensitivity analysis;
all results below use the fragility model.

\subsection{Population-Conserving Diurnal Exposure}
\label{sec:diurnal}
WorldPop is a nighttime residential surface. To estimate exposure at hour
$t$, occupancy-class multipliers $m_{o,t}$ act as \textit{relative weights}
and the citywide total is renormalized to the static baseline:
\begin{equation}
p_{i,t} = p_i \, m_{o(i),t} \cdot
\frac{\sum_j p_j}{\sum_j p_j \, m_{o(j),t}},
\end{equation}
so that people move between buildings but never leave the study area.
Expected casualties follow \textbf{HAZUS severity-stratified rates}
\cite{hazus} applied to collapsed-building occupants by construction class,
replacing the earlier flat 35\% injury-plus-fatality proxy. Occupancy
multipliers by period and land-use class follow HAZUS-inspired diurnal
patterns (residential dominance at night, commercial/education peaks at
midday) and are documented as \textit{scenario weights} in the released
configuration; perturbing non-residential multipliers by $\pm$20\% changes
midday casualties by less than 1\% under population conservation
(Section~\ref{sec:sensitivity}).

\subsection{Emergency Response}
\label{sec:response}
Each building is assigned its nearest health facility. Where the OpenStreetMap
road graph connects building and facility, travel time follows a
\textbf{network shortest path} (Dijkstra on the largest connected component);
otherwise a straight-line fallback is used. Only 5.9\% of buildings are
network-routed in the current graph --- the routing module is included as a
\textbf{proof-of-concept} for graph-complete OSM extraction, not as a
citywide travel-time model. Emergency capacity is estimated from OSM amenity
class (a known weakness; DGHS bed censuses would be the authoritative
replacement). A response-time proxy combines road-class travel speed
with local blockage delays that grow with the density of high-collapse-probability
buildings within 25~m of the route. Rescue priority weights normalized expected
casualties (0.65), collapse probability (0.20), and inverse response time (0.15).

\subsection{Lifeline Dependency Graph and Cascade Simulation}
\label{sec:cascade}
The lifeline graph contains 2{,}479 nodes --- 61 power nodes, 8
communication towers, 1{,}195 bridges, 728 hospitals, and 487 building
zones ($\sim$1.1~km aggregation cells carrying population and casualties)
--- and 9{,}737 dependency edges wired by spatial proximity: grid
interconnects between substations ($k=2$), power feeds to towers, hospitals,
and zones, communication links to hospitals and zones, and bridge access to
zones within 1.5~km.

Each of 300 Monte Carlo iterations proceeds in three steps. (1)~Direct
failures are sampled per infrastructure node from a logistic fragility on
amplified PGA (median capacities 0.35--0.50~g by node type) plus a
liquefaction boost. (2)~Failures propagate: a functioning substation fails
with probability $1-(1-0.30)^{d}$ when $d$ interconnected neighbours are
down (load redistribution \cite{motter2002}); towers and hospitals fail
with probability 0.65 and 0.30 respectively when all feeding substations
are lost \cite{ouyang2014} (backup generation folded into the conditional
probability; values are literature-informed and perturbed in sensitivity
analysis). (3)~Zones lose a service when all providers are down, and rescue
delay multiplies by 1.60 (communication lost), 1.25 (power lost), and 1.50
(bridge out), composing additively. We use $N = 300$ Monte Carlo iterations
after a convergence check showing population-weighted delay factor stabilizes
within $\pm$2\% of the $N = 500$ mean for $N \geq 200$
(\texttt{scripts/mc\_convergence\_study.py}). The ensemble runs six scenarios
(four Dauki magnitudes for training, two Madhupur magnitudes for
generalization testing).

\subsection{Physics-Informed Surrogate Emulator}
\label{sec:surrogate}
The cascade Monte Carlo simulation (Section~\ref{sec:cascade}) is the
authoritative physics model; a \textbf{physics-informed surrogate emulator}
accelerates it for interactive digital-twin queries. We implement the emulator
as a two-layer graph neural network trained on Monte Carlo labels. After
architecture ablation (Section~\ref{sec:sensitivity}), \textbf{GraphSAGE}
\cite{hamilton2017} is adopted for production inference (lowest test MAE,
$R^2 = 0.992$); GCN \cite{kipf2017} is retained as the literature-standard
spectral baseline in the comparison table. The production model uses a
symmetric normalized adjacency $\hat{A} = D^{-1/2}(A+I)D^{-1/2}$ over the
1{,}992 infrastructure nodes and predicts each node's \textit{total}
(direct plus cascaded) failure probability --- not fragility alone.
Twelve input features comprise node-type one-hot, amplified PGA, liquefaction
index, in/out degree, direct fragility probability, and scenario magnitude.
Two training regimes are reported:
(a)~75/25 node hold-out on the $M_w$~7.2 Dauki graph; and (b)~cross-scenario
training on four Dauki graphs with evaluation on all nodes from held-out
Madhupur scenarios. Training uses mean-squared-error loss and the Adam
optimizer (500 epochs, 32 hidden units for single-scenario; 48 for
cross-scenario). The comparison baseline is the physics-only direct
fragility, which cannot represent cascades by construction.

\subsection{Earthquake Resilience Index}
\label{sec:eri}
A composite Earthquake Resilience Index (ERI) is computed for each of the
487 building zones as a weighted sum of four normalized components:
structural resilience ($1 - \bar{P}_{\mathrm{collapse}}$, weight 0.30),
emergency-capacity headroom (weight 0.25), lifeline robustness
(inverse cascade delay factor, weight 0.25), and accessibility
(inverse normalized response time, weight 0.20). Weights reflect equal
stakeholder emphasis for exposition; one-at-a-time $\pm$0.10 perturbations
change citywide ERI by less than 2 points (Section~\ref{sec:sensitivity}).
Citywide ERI is the population-weighted mean; zones are tiered low ($<$40),
moderate (40--55), high (55--70), and very high ($\geq$70).

\section{Experiments}
\label{sec:experiments}

\subsection{Experimental Setup}
All analyses were implemented in Python (GeoPandas, NetworkX, PyTorch,
XGBoost); configuration lives in versioned YAML files and every figure and
table below is regenerated by a single scripted pipeline (see
Section~\ref{sec:code_availability}).

\subsection{Hazard and Building Vulnerability}
Bedrock PGA over the study area is 0.145--0.178~g (mean 0.161~g),
soil-amplified PGA 0.188--0.252~g (mean 0.224~g), and mean MMI 8.08.
Figure~\ref{fig:hazard} maps amplified PGA; amplification concentrates in
the eastern and riverine soft-soil zones.

\casinlinefig[0.9\linewidth]{F5_2_amplified_pga.png}{Soil-amplified peak ground acceleration for the $M_w$~7.2 Dauki scenario.}{fig:hazard}

Of 706{,}115 buildings, mean collapse probability is 0.40 (median 0.40),
with only 0.33\% exceeding a 0.5 threshold (2{,}323 buildings) under the
logic-tree GMPE and Vs30 amplification --- most stock falls in the moderate
band (0.2--0.5). Fragility ordering across construction classes remains
physically sensible (Table~\ref{tab:construction}, Figures~\ref{fig:collapse}
and \ref{fig:byclass}): informal construction is most fragile (mean 0.657)
and engineered reinforced concrete least (0.141). Notably, 95.4\% of buildings
carry no OSM construction tag and default to the ``unknown'' fragility
class --- the dominant epistemic uncertainty in the vulnerability layer.

\begin{table}[pos=htbp]
  \centering
  \caption{Collapse probability by construction class ($M_w$~7.2 Dauki).}
  \label{tab:construction}
  \small
  \begin{tabular}{lrr}
    \toprule
    Construction class & Buildings & Mean $P_{\mathrm{collapse}}$ \\
    \midrule
    Informal            & 176       & 0.657 \\
    Unknown (untagged)  & 673{,}392 & 0.401 \\
    Residential (tagged)& 26{,}374  & 0.359 \\
    Industrial          & 415       & 0.267 \\
    Institutional       & 890       & 0.223 \\
    Commercial          & 1{,}553   & 0.200 \\
    Reinforced concrete & 3{,}315   & 0.141 \\
    \bottomrule
  \end{tabular}
\end{table}

\casinlinefig[0.9\linewidth]{F6_1_collapse_probability_map.png}{Building-level collapse probability across Greater Dhaka.}{fig:collapse}

\casinlinefig[0.85\linewidth]{F6_2_collapse_by_construction_type.pdf}{Mean collapse probability by normalized construction class.}{fig:byclass}

\subsection{Diurnal Exposure}
Under population conservation, total exposure is constant by construction;
what varies is its location (Table~\ref{tab:diurnal},
Figure~\ref{fig:diurnal}). Only $\approx$35{,}000 people occupy buildings
with collapse probability $\geq 0.5$ at any hour under the upgraded hazard
model; diurnal variation in high-risk exposure is $\approx$1\%. HAZUS-expected
casualties at midday total 322{,}837 and vary by less than 2\% across the
day. The flatness reflects both OSM occupancy incompleteness and the fact
that most buildings now fall below the 0.5 collapse threshold.

\begin{table}[pos=htbp]
  \centering
  \caption{Diurnal exposure ($M_w$~7.2 Dauki; population-conserving).}
  \label{tab:diurnal}
  \small
  \setlength{\tabcolsep}{4pt}
  \begin{tabular}{lrrr}
    \toprule
    Period & Population & High-risk bldgs & Expected casualties \\
    \midrule
    00:00 midnight & 4{,}475{,}588 & 35{,}038 & 328{,}634 \\
    08:00 commute  & 4{,}586{,}032 & 35{,}093 & 326{,}843 \\
    13:00 midday   & 4{,}593{,}956 & 35{,}210 & 322{,}837 \\
    18:00 evening  & 4{,}542{,}278 & 34{,}903 & 327{,}568 \\
    \bottomrule
  \end{tabular}
\end{table}

\casinlinefig[\linewidth]{F7_1_diurnal_population_casualties.pdf}{High-risk exposure and expected severe impacts by time of day.}{fig:diurnal}

\subsection{Emergency Response}
At midday, HAZUS-expected casualty demand (322{,}837) exceeds the
estimated emergency capacity of 243 of 728 mapped health facilities (33\%),
with the worst-affected facility receiving 38$\times$ its proxy capacity;
the mean response-time proxy is 7.2 minutes (network-routed where connected).
Rescue priority (Figure~\ref{fig:rescue}) concentrates on large-footprint,
high-occupancy structures rather than the single most fragile buildings,
reflecting the casualty-weighted score.

\casinlinefig[0.85\linewidth]{F8_2_rescue_priority_map.png}{Building-level rescue priority (midday scenario).}{fig:rescue}

\subsection{Cascading Lifeline Failure}
Cascades transform the lifeline picture (Table~\ref{tab:cascade},
Figure~\ref{fig:cascade}). Direct seismic fragility alone predicts mean
failure probabilities of 0.07--0.14 across node types; with grid
interdependency, substation failure rises to 0.600 (4.3$\times$),
communication-tower loss to 0.766 (6.0$\times$), and hospital functional
loss to 0.604 (7.1$\times$). Bridges, which depend on nothing upstream, are
unchanged --- a useful internal control. Population-weighted, residents face
a 67\% probability of power loss and 65\% of bridge-access disruption, and
the expected rescue-delay factor is 1.90 (95\% CI 1.88--1.93 across MC
iterations at $N = 300$).

The six-scenario ensemble reveals that \textbf{fault proximity dominates
magnitude}: a $M_w$~7.0 Madhupur rupture (closer to Dhaka) produces
near-total lifeline failure (100\% communication loss, delay factor 2.34)
whereas a $M_w$~7.2 Dauki event yields 68\% communication loss and delay
1.90 (Table~\ref{tab:ensemble}).

\begin{table}[pos=htbp]
  \centering
  \caption{Population-weighted cascade impact by scenario (300 MC iterations).}
  \label{tab:ensemble}
  \small
  \setlength{\tabcolsep}{3pt}
  \begin{tabular}{lccc}
    \toprule
    Scenario & $P$(power lost) & $P$(comm.\ lost) & Delay factor \\
    \midrule
    $M_w$~6.8 Dauki   & 0.47 & 0.47 & 1.65$\times$ \\
    $M_w$~7.0 Dauki   & 0.59 & 0.58 & 1.78$\times$ \\
    $M_w$~7.2 Dauki   & 0.67 & 0.68 & 1.90$\times$ \\
    $M_w$~7.5 Dauki   & 0.83 & 0.85 & 2.09$\times$ \\
    $M_w$~6.5 Madhupur & 0.98 & 0.99 & 2.31$\times$ \\
    $M_w$~7.0 Madhupur & 1.00 & 1.00 & 2.34$\times$ \\
    \bottomrule
  \end{tabular}
\end{table}

\begin{table}[pos=htbp]
  \centering
  \caption{Mean failure probability by node type ($M_w$~7.2 Dauki; 300 MC iterations).}
  \label{tab:cascade}
  \small
  \setlength{\tabcolsep}{4pt}
  \begin{tabular}{lccc}
    \toprule
    Node type & Direct & Total (cascade) & Amplification \\
    \midrule
    Power (61)       & 0.139 & 0.600 & 4.3$\times$ \\
    Comm.\ tower (8) & 0.127 & 0.766 & 6.0$\times$ \\
    Hospital (728)   & 0.085 & 0.604 & 7.1$\times$ \\
    Bridge (1{,}195) & 0.069 & 0.068 & 1.0$\times$ \\
    \bottomrule
  \end{tabular}
\end{table}

\casinlinefig[0.85\linewidth]{F9_1_cascade_amplification.pdf}{Cascade amplification of failure probability by node type.}{fig:cascade}

\subsection{Surrogate Emulator Performance}
\label{sec:surrogate_results}
The physics-informed surrogate emulator is evaluated against its Monte Carlo
training labels. On a 75/25 node hold-out of the $M_w$~7.2 Dauki graph, the
production GraphSAGE emulator
predicts total failure probability with mean absolute error 0.016 and
$R^2 = 0.992$ (Figure~\ref{fig:gnn}); the physics-only baseline achieves
mean absolute error 0.231 and $R^2 = -0.55$. Cross-scenario training on
four Dauki graphs and evaluation on held-out Madhupur scenarios yields
$R^2 = 0.883$ (MAE 0.078) versus a direct-fragility baseline of
$R^2 = 0.305$ (Table~\ref{tab:gnn}), confirming transferable cascade
structure rather than single-graph memorization. Inference over the full
graph takes milliseconds, versus minutes for Monte Carlo simulation.

\begin{table}[pos=htbp]
  \centering
  \caption{Surrogate emulator quality: single-scenario node hold-out vs.\ cross-scenario generalization.}
  \label{tab:gnn}
  \small
  \begin{tabular}{lcc}
    \toprule
    Model & MAE & $R^2$ \\
    \midrule
    GraphSAGE ($M_w$~7.2 node hold-out)     & 0.016 & 0.992 \\
    GraphSAGE (held-out Madhupur scenarios) & 0.078 & 0.883 \\
    Direct fragility (Madhupur test)  & 0.117 & 0.305 \\
    \bottomrule
  \end{tabular}
\end{table}

\casinlinefig[0.75\linewidth]{F9_4_gnn_vs_montecarlo.pdf}{Physics-informed surrogate emulator versus Monte Carlo total failure probability on held-out nodes.}{fig:gnn}

\subsection{Earthquake Resilience Index}
Citywide ERI is 70.9 (very-high tier; mean zone ERI 69.4). Of 487 zones,
252 fall in the very-high tier ($\geq$70) and 235 in the high tier (55--70).
Component scores (0--1, higher = more resilient) are structural 0.60,
emergency capacity 0.86, lifeline robustness 0.61, and accessibility 0.81.
Lifeline and structural components are the binding constraints under the
reference scenario; emergency-capacity headroom is higher once HAZUS
severity rates replace the earlier flat casualty proxy. Manual YAML weights
(0.30/0.25/0.25/0.20) are compared with an entropy--hybrid scheme that
assigns data-driven weights to spatially varying components and retains the
manual emergency weight where zonal hospital load is uniform; citywide ERI
shifts from 70.9 to 76.1 under the hybrid scheme (Section~\ref{sec:sensitivity}).

\subsection{External Validation}
\label{sec:external_validation}
Internal checks (literature PGA/MMI envelopes at 158~km) are supplemented by
distance--magnitude envelope tests and a historical analog replay
(\texttt{scripts/run\_external\_validation.py}). At 150~km and $M_w$~7.2 the
logic-tree bedrock PGA (0.169~g) falls within the CDMP-style
0.06--0.18~g envelope used for Dhaka scenario classes; at 100~km and
200~km the model sits at the upper and lower bounds respectively, reflecting
GMPE conservatism at short range and attenuation at long range. Replaying the
1897 Great Assam earthquake ($M_w \approx 8.1$, 250~km from Dhaka) yields
model MMI 8.3 versus macroseismic reports of VI--VII ($\approx$6.5--7.5 MMI)
\cite{oldham1899} --- the GMPE replay \textit{overpredicts} shaking, so it is
used as an order-of-magnitude consistency check, not damage calibration.
ShakeMap raster comparison, observed hospital disruption, and utility outage
validation are not yet possible with open Bangladesh feeds; these remain the
principal external-validation gaps (Section~\ref{sec:limitations}).

\subsection{Sensitivity Analysis and Uncertainty}
\label{sec:sensitivity}
Because many sub-models rely on open-data proxies, we report targeted
one-at-a-time sensitivity studies (scripts in \texttt{scripts/run\_sensitivity\_studies.py},
\texttt{scripts/gnn\_architecture\_ablation.py},
\texttt{scripts/mc\_convergence\_study.py}) rather than only point estimates.
(i)~\textbf{GMPE logic-tree weights:} 40/60 and 60/40 branch splits change
mean amplified PGA by $\leq$3.6\%. (ii)~\textbf{Occupancy multipliers:}
$\pm$20\% on non-residential classes (residential fixed, totals conserved)
changes midday casualties by 0.9\%. (iii)~\textbf{Cascade conditional
probabilities:} $\pm$0.15 perturbations change the citywide delay factor by
up to 11.8\% (substation overload rate most influential). (iv)~\textbf{ERI
weights:} $\pm$0.10 one-at-a-time perturbations change citywide ERI by
$\leq$1.7 points. (v)~\textbf{Monte Carlo count:} population-weighted delay
factor varies by $<$1.2\% between $N = 50$ and $N = 500$; $N = 300$ differs
from $N = 500$ by 0.4\%. (vi)~\textbf{Surrogate architecture:} GraphSAGE is
adopted for production inference (lowest test MAE 0.016, $R^2 = 0.992$); GCN,
GAT, and MLP are within 0.004 MAE; topology-free MLP underperforms graph models
(Table~\ref{tab:gnn_ablation}). (vii)~\textbf{Vs30:} $\pm$20\% perturbation
changes mean $P_{\mathrm{collapse}}$ by up to 8.6\%. (viii)~\textbf{ERI
weighting:} entropy--hybrid weights (spatially varying components only) shift
citywide ERI by 5.2 points versus manual weights (70.9 vs.\ 76.1).

\textbf{Latin Hypercube uncertainty propagation} ($N = 128$; six parameters:
GMPE weight, fragility median shift, logistic slope, cascade conditional
probabilities, non-residential occupancy) jointly samples epistemic uncertainty
(\texttt{scripts/run\_latin\_hypercube\_uq.py}). Table~\ref{tab:uq_ci}
reports 95\% intervals on citywide metrics; casualties and mean collapse
probability vary substantially under coupled perturbations, whereas
one-at-a-time tests suggested smaller marginal effects.

\begin{table}[pos=htbp]
  \centering
  \caption{Latin Hypercube 95\% confidence intervals ($N = 128$; $M_w$~7.2 Dauki midday).}
  \label{tab:uq_ci}
  \small
  \setlength{\tabcolsep}{3pt}
  \begin{tabular}{lcc}
    \toprule
    Metric & Mean & 95\% CI \\
    \midrule
    Mean $P_{\mathrm{collapse}}$ & 0.41 & 0.33--0.47 \\
    Midday casualties (millions) & 0.32 & 0.19--0.41 \\
    Cascade delay factor & 1.90 & 1.60--2.22 \\
    Citywide ERI & 69.1 & 64.0--75.2 \\
    \bottomrule
  \end{tabular}
\end{table}

\begin{table}[pos=htbp]
  \centering
  \caption{Surrogate emulator architecture ablation (75/25 node hold-out; $M_w$~7.2 Dauki).}
  \label{tab:gnn_ablation}
  \small
  \begin{tabular}{lcc}
    \toprule
    Architecture & MAE & $R^2$ \\
    \midrule
    GraphSAGE       & 0.016 & 0.992 \\
    GAT (gated blend) & 0.017 & 0.991 \\
    GCN             & 0.018 & 0.988 \\
    MLP (no graph)  & 0.022 & 0.976 \\
    GIN             & 0.052 & 0.950 \\
    Direct fragility & 0.222 & $-0.56$ \\
    \bottomrule
  \end{tabular}
\end{table}

\section{Discussion}
Compared with ward-level Dhaka assessments \cite{cdmp2009,islam2020}, our
bedrock PGA (mean 0.161~g) sits within the CDMP envelope for a large eastern
fault event at 150--200~km, while building-resolved fragility exposes
moderate-band dominance (mean $P_{\mathrm{collapse}} = 0.40$) rather than
the binary high/low splits typical of index-based studies. HAZUS-expected
casualties ($\approx$0.32~M at midday) are planning-order estimates --- lower
than naive 35\% flat-rate proxies but not directly comparable to historical
fatality counts from distant analogues (e.g., 1897 Assam) without event
matching. Relative to OpenQuake/GEM portfolio tools \cite{pagani2014,gem2023},
QuakeTwin's contribution is the \textit{coupled} city-scale chain from OSM
footprints through diurnal exposure, routing, cascade, and surrogate emulation.

Three findings carry direct policy weight. First, the vulnerability layer is
dominated by untagged, presumably non-engineered stock in the eastern and
riverine soft-soil zones where amplification and liquefaction coincide ---
the same neighbourhoods that rank worst on cascade-driven rescue delay.
Retrofitting and code-enforcement effort should be prioritized there, and a
structural tagging campaign (even a coarse engineered/non-engineered
dichotomy) would collapse the largest single uncertainty in this study.

Second, HAZUS-expected casualty demand exceeding one-third of mapped
emergency capacity implies that nearest-hospital routing is not a viable
response doctrine for this scenario class; pre-planned casualty distribution
and field-hospital pre-positioning are required.

Third, the 4--7$\times$ cascade amplification of lifeline failure means
hardening a small number of substations and guaranteeing hospital backup
power buys disproportionate system resilience --- these are cheaper
interventions than structural retrofitting at scale, and the dependency
graph identifies exactly which nodes matter most. The Madhupur ensemble
shows that fault proximity can matter more than magnitude: planners must
stress-test multiple epicenters, not a single reference rupture.

The cross-scenario surrogate emulator ($R^2 = 0.883$ on held-out Madhupur graphs) converts
minutes-scale Monte Carlo simulation into a milliseconds-scale function with modest
accuracy loss relative to single-scenario training ($R^2 = 0.992$). This
trade-off is acceptable for interactive epicenter sweeps in the digital twin.
Architecturally, QuakeTwin separates \textit{state} (building and infrastructure
graphs), \textit{simulation} (hazard, fragility, cascade Monte Carlo), and
\textit{surrogate inference} (graph-based emulator), enabling a future Phase~6 dashboard where
seismic feeds or hospital status updates trigger re-evaluation in milliseconds
via the surrogate while retaining the physics pipeline for calibration runs.

\section{Limitations}
\label{sec:limitations}
All quantitative claims are conditional on open-data completeness and stated
assumptions. First, 95\% of buildings lack OSM construction and occupancy
tags, defaulting to class medians; the diurnal flatness partly reflects this
incompleteness. Second, soil zones combine Vs30 sampling with
literature-informed proxies pending borehole validation. Third, network
shortest-path routing connects only 5.9\% of buildings to hospitals in the
current road graph; most assignments still use straight-line fallbacks.
Fourth, hospital capacities derive from amenity tags. Fifth, only eight
communication towers are mapped, dependency wiring is proximity-based rather
than actual feeder or backhaul topology, and conditional failure
probabilities are literature-informed assumptions. Sixth, cross-scenario surrogate
$R^2$ (0.883) is lower than single-scenario performance (0.992) --- both
should be reported. Seventh, external validation is limited to literature
PGA/MMI envelopes and a qualitative 1897 Assam replay; there is no ShakeMap
raster comparison, post-event damage survey, or observed hospital/utility
outage calibration. The pipeline is designed so that each proxy can be
swapped for authoritative data (borehole logs, DPDC/DESCO feeders, DGHS bed
censuses) without structural change.

\section{Conclusion}
QuakeTwin demonstrates that a building-resolved, temporally dynamic,
cascade-aware earthquake digital twin of a data-scarce megacity can be
assembled entirely from open data and open source. For a $M_w$~7.2 Dauki
scenario, the twin estimates mean collapse probability 0.40 (0.33\% of
buildings exceed 0.5), HAZUS-expected casualties of roughly 0.32~million
at midday with less than 2\% diurnal variation, casualty demand overwhelming
one-third of mapped health-facility capacity, and cascading lifeline failure
--- invisible to direct fragility analysis --- raising expected rescue delay
by a factor of 1.90 (95\% CI 1.60--2.22 under Latin Hypercube propagation).
A nearer Madhupur rupture produces worse lifeline
outcomes despite lower magnitude. Latin Hypercube analysis yields 95\% CIs of
0.33--0.47 on mean collapse probability and 0.19--0.41~million on midday
casualties under jointly sampled epistemic uncertainty.

Principal findings include: (i) GMPE logic-tree and Vs30 amplification yield
defensible moderate-band fragility even under heavy tag incompleteness;
(ii) population conservation is essential for defensible diurnal exposure
claims; (iii) direct fragility understates lifeline failure by 4--7$\times$
once grid interdependency is modelled; (iv) fault proximity can dominate
magnitude in cascade outcomes; (v) physics-informed surrogate emulators generalize
across scenarios with $R^2 = 0.883$ at millisecond latency; and (vi) the composite ERI (70.9
citywide) localizes where structural and lifeline interventions matter most.
Future work extends the twin with a CesiumJS three-dimensional front end,
full road-graph connectivity for routing, Open Buildings height integration,
and authoritative utility and geotechnical data as they become available.

\declarationspage

\bibliography{cas-refs}

\end{document}
